\documentclass[10pt,a4paper]{amsart}
\usepackage{amsmath,amssymb,amsthm,graphicx}
\usepackage{verbatim}
\newenvironment{reference}{\comment}{\endcomment}
%\newenvironment{reference}{}{}
\newenvironment{slogan}{\comment}{\endcomment}
\newenvironment{history}{\comment}{\endcomment}

% For commutative diagrams we use Xy-pic
\usepackage[all]{xy}

% We use 2cell for 2-commutative diagrams.
\xyoption{2cell}
% Theorem environments.
%
\theoremstyle{plain}
\newtheorem{theorem}[subsection]{Theorem}
\newtheorem{proposition}[subsection]{Proposition}
\newtheorem{lemma}[subsection]{Lemma}

\theoremstyle{definition}
\newtheorem{definition}[subsection]{Definition}
\newtheorem{example}[subsection]{Example}
\newtheorem{exercise}[subsection]{Exercise}
\newtheorem{situation}[subsection]{Situation}

\theoremstyle{remark}
\newtheorem{remark}[subsection]{Remark}
\newtheorem{remarks}[subsection]{Remarks}

\numberwithin{equation}{subsection}

% Macros
%
\def\lim{\mathop{\mathrm{lim}}\nolimits}
\def\colim{\mathop{\mathrm{colim}}\nolimits}
\def\Spec{\mathop{\mathrm{Spec}}}
\def\Hom{\mathop{\mathrm{Hom}}\nolimits}
\def\Ext{\mathop{\mathrm{Ext}}\nolimits}
\def\SheafHom{\mathop{\mathcal{H}\!\mathit{om}}\nolimits}
\def\SheafExt{\mathop{\mathcal{E}\!\mathit{xt}}\nolimits}
\def\Sch{\mathit{Sch}}
\def\Mor{\mathop{\mathrm{Mor}}\nolimits}
\def\Ob{\mathop{\mathrm{Ob}}\nolimits}
\def\Sh{\mathop{\mathit{Sh}}\nolimits}
\def\NL{\mathop{N\!L}\nolimits}
\def\CH{\mathop{\mathrm{CH}}\nolimits}
\def\proetale{{pro\text{-}\acute{e}tale}}
\def\etale{{\acute{e}tale}}
\def\QCoh{\mathit{QCoh}}
\def\Ker{\mathop{\mathrm{Ker}}}
\def\Im{\mathop{\mathrm{Im}}}
\def\Coker{\mathop{\mathrm{Coker}}}
\def\Coim{\mathop{\mathrm{Coim}}}

% Boxtimes
%
\DeclareMathSymbol{\boxtimes}{\mathbin}{AMSa}{"02}

%
% Macros for moduli stacks/spaces
%
\def\QCohstack{\mathcal{QC}\!\mathit{oh}}
\def\Cohstack{\mathcal{C}\!\mathit{oh}}
\def\Spacesstack{\mathcal{S}\!\mathit{paces}}
\def\Quotfunctor{\mathrm{Quot}}
\def\Hilbfunctor{\mathrm{Hilb}}
\def\Curvesstack{\mathcal{C}\!\mathit{urves}}
\def\Polarizedstack{\mathcal{P}\!\mathit{olarized}}
\def\Complexesstack{\mathcal{C}\!\mathit{omplexes}}
% \Pic is the operator that assigns to X its picard group, usage \Pic(X)
% \Picardstack_{X/B} denotes the Picard stack of X over B
% \Picardfunctor_{X/B} denotes the Picard functor of X over B
\def\Pic{\mathop{\mathrm{Pic}}\nolimits}
\def\Picardstack{\mathcal{P}\!\mathit{ic}}
\def\Picardfunctor{\mathrm{Pic}}
\def\Deformationcategory{\mathcal{D}\!\mathit{ef}}
\usepackage{cleveref}
\crefname{originallemma}{original lemma}{original lemmas}
\Crefname{originallemma}{Original lemma}{Original lemmas}
\makeatletter
\newlabel{divisors-lemma-blowup-fitting-ideal}{{0ESN}{source}{}{sourcepointer-1}{}}
\newlabel{divisors-lemma-blowup-fitting-ideal@cref}{{[originallemma][0][]0ESN}{[1][0][]source}{}{}{}}
\makeatother
\crefname{originallemma}{original lemma}{original lemmas}
\Crefname{originallemma}{Original lemma}{Original lemmas}
\makeatletter
\newlabel{duality-lemma-dualizing-module}{{0AWK}{source}{}{sourcepointer-2}{}}
\newlabel{duality-lemma-dualizing-module@cref}{{[originallemma][0][]0AWK}{[1][0][]source}{}{}{}}
\makeatother
\crefname{originallemma}{original lemma}{original lemmas}
\Crefname{originallemma}{Original lemma}{Original lemmas}
\makeatletter
\newlabel{dualizing-lemma-regular-gorenstein}{{0AWX}{source}{}{sourcepointer-3}{}}
\newlabel{dualizing-lemma-regular-gorenstein@cref}{{[originallemma][0][]0AWX}{[1][0][]source}{}{}{}}
\makeatother
\crefname{originalreference}{original reference}{original references}
\Crefname{originalreference}{Original reference}{Original references}
\makeatletter
\newlabel{lemma-blow-up-normal-rational}{{0B4Z}{source}{}{sourcepointer-4}{}}
\newlabel{lemma-blow-up-normal-rational@cref}{{[originalreference][0][]0B4Z}{[1][0][]source}{}{}{}}
\makeatother
\crefname{originalreference}{original reference}{original references}
\Crefname{originalreference}{Original reference}{Original references}
\makeatletter
\newlabel{lemma-blowup-regular}{{0AGR}{source}{}{sourcepointer-5}{}}
\newlabel{lemma-blowup-regular@cref}{{[originalreference][0][]0AGR}{[1][0][]source}{}{}{}}
\makeatother
\crefname{originalreference}{original reference}{original references}
\Crefname{originalreference}{Original reference}{Original references}
\makeatletter
\newlabel{lemma-dominate-by-normalized-blowing-up}{{0BBT}{source}{}{sourcepointer-6}{}}
\newlabel{lemma-dominate-by-normalized-blowing-up@cref}{{[originalreference][0][]0BBT}{[1][0][]source}{}{}{}}
\makeatother
\crefname{originalreference}{original reference}{original references}
\Crefname{originalreference}{Original reference}{Original references}
\makeatletter
\newlabel{lemma-dualizing-blow-up-rational}{{0B64}{source}{}{sourcepointer-7}{}}
\newlabel{lemma-dualizing-blow-up-rational@cref}{{[originalreference][0][]0B64}{[1][0][]source}{}{}{}}
\makeatother
\crefname{originalreference}{original reference}{original references}
\Crefname{originalreference}{Original reference}{Original references}
\makeatletter
\newlabel{lemma-dualizing-rational}{{0BBU}{source}{}{sourcepointer-8}{}}
\newlabel{lemma-dualizing-rational@cref}{{[originalreference][0][]0BBU}{[1][0][]source}{}{}{}}
\makeatother
\title{A rational surface singularity with a dualizing complex is made Gorenstein by blowups in singular closed points}
\author{The Stacks Project Authors}
\date{stacks-project:resolve.tex:local-snapshot-20261012; GNU FDL 1.2 or later}
\begin{document}
\maketitle
\noindent\textit{Editorial note.} All delivered mathematical text is the author's own, copied from the pinned Stacks Project chapter file resolve.tex; no statement or proof step was rewritten, reordered or removed.\par\medskip
\noindent\textit{Editorial note.} The counted proof is the author's complete argument: the construction of a blowup making the strict transform of the dualizing module invertible, its domination by a sequence of normalized blowups at singular closed points with the normality of every term, and the pointwise verification of invertibility at regular and singular points closed by the torsion freeness of the dualizing module.\par\medskip
\noindent\textit{Editorial note.} The label lemma-rational-to-gorenstein names the counted statement, and the proof environment immediately following it in the source is delivered verbatim as its proof.\par\medskip
\noindent\textit{Editorial note.} The counted statement is carried by the author's own lemma environment in the chapter source.\par\medskip
\noindent\textit{Editorial note.} Pre-open gate: python operations/agent-authoring/tools/precard\_gate.py --source resolve.tex --label lemma-rational-to-gorenstein --start 2687 --end 2743 --local-id 28081 returned RESULT=OK (L28081 free, occupancy FREE, public interval clear). Pre-assembly tag self-check: canonical tag resolved and not present in reference/public-items.jsonl.\par\medskip
\noindent\textit{Editorial note.} The statement and proof are copied from the pinned chapter file; the cited inputs are kept as external pointers with their canonical Stacks tags.\par\medskip
\noindent\textit{Allowed standard prerequisites (editorial).}\begin{itemize}
\item Invertibility of the strict transform of the dualizing module. Divisors, the lemma producing a blowup with respect to the Fitting ideal such that the strict transform of a given coherent module is invertible. Supplies the blowup over which the strict transform of the dualizing module is an invertible sheaf.
\item Domination by normalized blowups and normality of the blowups. Same chapter, the lemma dominating a proper birational morphism by normalized blowups at closed points, and the lemma that the blowup of the maximal ideal of a normal rational surface singularity is normal. Supply the sequence of normalized blowups dominating the chosen blowup together with the fact that each step is an ordinary blowup and every term is normal.
\item Dualizing modules of regular and of blown up rational singularities, and injectivity of the comparison map. Dualizing Complexes, the lemma that a regular scheme is Gorenstein, and Same chapter, the lemmas describing the dualizing module of the blowup and the surjectivity of the comparison map along such a blowup. Supply the invertibility at regular points and the surjectivity on stalks along the sequence at singular points.
\item Torsion freeness of the dualizing module. Duality for Schemes, the lemma that the dualizing module of a normal scheme is torsion free. Supplies the final identification of the quotient by torsion with the dualizing module of the last term.
\end{itemize}\medskip
\ifdefined\hypertarget\hypertarget{sourcepointer-1}{}\fi
\noindent\textit{Source reference.} Original lemma 0ESN: sources/stacks-chow/tags.txt, source lines 18488--18488. Prerequisite pointer invoked at proof line(s) 2708; it names a standard fact of a sibling chapter that the counted argument uses, and every step of the counted proof built on it is retained verbatim in the delivered body.\par\medskip
\ifdefined\hypertarget\hypertarget{sourcepointer-2}{}\fi
\noindent\textit{Source reference.} Original lemma 0AWK: sources/stacks-chow/tags.txt, source lines 13723--13723. Prerequisite pointer invoked at proof line(s) 2740; it names a standard fact of a sibling chapter that the counted argument uses, and every step of the counted proof built on it is retained verbatim in the delivered body.\par\medskip
\ifdefined\hypertarget\hypertarget{sourcepointer-3}{}\fi
\noindent\textit{Source reference.} Original lemma 0AWX: sources/stacks-chow/tags.txt, source lines 13735--13735. Prerequisite pointer invoked at proof line(s) 2723; it names a standard fact of a sibling chapter that the counted argument uses, and every step of the counted proof built on it is retained verbatim in the delivered body.\par\medskip
\ifdefined\hypertarget\hypertarget{sourcepointer-4}{}\fi
\noindent\textit{Source reference.} Original reference 0B4Z: sources/stacks-resolve/pinned-original-source.tex, source lines 2534--2534. Prerequisite pointer invoked at proof line(s) 2714; it names a standard fact of the same chapter that the counted argument uses, and every step of the counted proof built on it is retained verbatim in the delivered body.\par\medskip
\ifdefined\hypertarget\hypertarget{sourcepointer-5}{}\fi
\noindent\textit{Source reference.} Original reference 0AGR: sources/stacks-resolve/pinned-original-source.tex, source lines 727--727. Prerequisite pointer invoked at proof line(s) 2726; it names a standard fact of the same chapter that the counted argument uses, and every step of the counted proof built on it is retained verbatim in the delivered body.\par\medskip
\ifdefined\hypertarget\hypertarget{sourcepointer-6}{}\fi
\noindent\textit{Source reference.} Original reference 0BBT: sources/stacks-resolve/pinned-original-source.tex, source lines 1277--1277. Prerequisite pointer invoked at proof line(s) 2709; it names a standard fact of the same chapter that the counted argument uses, and every step of the counted proof built on it is retained verbatim in the delivered body.\par\medskip
\ifdefined\hypertarget\hypertarget{sourcepointer-7}{}\fi
\noindent\textit{Source reference.} Original reference 0B64: sources/stacks-resolve/pinned-original-source.tex, source lines 2641--2641. Prerequisite pointer invoked at proof line(s) 2731; it names a standard fact of the same chapter that the counted argument uses, and every step of the counted proof built on it is retained verbatim in the delivered body.\par\medskip
\ifdefined\hypertarget\hypertarget{sourcepointer-8}{}\fi
\noindent\textit{Source reference.} Original reference 0BBU: sources/stacks-resolve/pinned-original-source.tex, source lines 2618--2618. Prerequisite pointer invoked at proof line(s) 2728; it names a standard fact of the same chapter that the counted argument uses, and every step of the counted proof built on it is retained verbatim in the delivered body.\par\medskip
\begin{lemma}
\label{lemma-rational-to-gorenstein}
Let $(A, \mathfrak m, \kappa)$ be a local normal Nagata domain of
dimension $2$ which defines a rational singularity. Assume $A$ has
a dualizing complex. Then there exists a finite sequence of blowups in
singular closed points
$$
X = X_n \to X_{n - 1} \to \ldots \to X_1 \to X_0 = \Spec(A)
$$
such that $X_i$ is normal for each $i$ and such that
the dualizing sheaf $\omega_X$ of $X$ is an invertible
$\mathcal{O}_X$-module.
\end{lemma}

\begin{proof}
The dualizing module $\omega_A$ is a finite $A$-module whose stalk at
the generic point is invertible. Namely, $\omega_A \otimes_A K$
is a dualizing module for the fraction field $K$ of $A$, hence has
rank $1$. Thus there exists a blowup $b : Y \to \Spec(A)$ such that
the strict transform of $\omega_A$ with respect to $b$ is an invertible
$\mathcal{O}_Y$-module, see
Divisors, Lemma \ref{divisors-lemma-blowup-fitting-ideal}.
By Lemma \ref{lemma-dominate-by-normalized-blowing-up}
we can choose a sequence of normalized blowups
$$
X_n \to X_{n - 1} \to \ldots \to X_1 \to \Spec(A)
$$
such that $X_n$ dominates $Y$. By Lemma \ref{lemma-blow-up-normal-rational}
and arguing by induction each $X_i \to X_{i - 1}$ is simply a blowing up.

\medskip\noindent
We claim that $\omega_{X_n}$ is invertible. Since $\omega_{X_n}$
is a coherent $\mathcal{O}_{X_n}$-module, it suffices to see its stalks
are invertible modules. If $x \in X_n$ is a regular point, then this is
clear from the fact that regular schemes are
Gorenstein (Dualizing Complexes, Lemma
\ref{dualizing-lemma-regular-gorenstein}). If $x$ is a singular point of
$X_n$, then each of the images $x_i \in X_i$ of $x$ is a singular point
(because the blowup of a regular point is regular by
Lemma \ref{lemma-blowup-regular}).
Consider the canonical map $f_n^*\omega_A \to \omega_{X_n}$ of
Lemma \ref{lemma-dualizing-rational}. For each $i$ the morphism
$X_{i + 1} \to X_i$ is either a blowup of $x_i$ or an isomorphism
at $x_i$. Since $x_i$ is always a singular point, it follows from
Lemma \ref{lemma-dualizing-blow-up-rational}
and induction that the maps $f_i^*\omega_A \to \omega_{X_i}$
is always surjective on stalks at $x_i$. Hence
$$
(f_n^*\omega_A)_x \longrightarrow \omega_{X_n, x}
$$
is surjective. On the other hand, by our choice of $b$ the quotient
of $f_n^*\omega_A$ by its torsion submodule is an invertible module
$\mathcal{L}$. Moreover, the dualizing module is torsion free
(Duality for Schemes, Lemma \ref{duality-lemma-dualizing-module}).
It follows that $\mathcal{L}_x \cong  \omega_{X_n, x}$ and the proof is
complete.
\end{proof}

\end{document}
